% GENERATED FILE. Edit canonical lessons, then run npm run generate:books.
% canonical-source-hash: fnv1a64:13c038e095e0c48e
% canonical-lessons: SA-C242-sulabhah, SA-C242-durlabhah, SA-C242-jatilah, SA-C242-sukarah, SA-C242-duskarah

\chapter{Easy to get, Hard to get, Complicated, Easy to do, Hard to do}
\label{ch:sa-easy-to-get-hard-to-get-complicated-easy-to-do-hard-to-do}
\hlchaptermodality{242}

\begin{chapteropening}
\textbf{By the end of this chapter:} \emph{I can say easy to get, hard to get, complicated, easy to do and hard to do.}

Complete the last lesson of chapter 242: I can say easy to get, hard to get, complicated, easy to do and hard to do.
\end{chapteropening}

\section[\texorpdfstring{\sk{सुलभः} (sulabhaḥ)}{सुलभः (sulabhaḥ)}]{\sk{सुलभः} (sulabhaḥ) — easy to get, cheap}

\label{lesson:SA-C242-sulabhah}



\begin{quote}
Before the new one: say the Sanskrit for that, then the Sanskrit for however.
\end{quote}

\subsection*{You'll want to know: \sk{सुलभः}}
\textbf{\sk{सुलभः}} — \emph{sulabhaḥ} — \textquotedblleft{}easy to get, cheap\textquotedblright{}.

\begin{grammarlens}[title={What you've built}]
One more word for this chapter.
\end{grammarlens}

\subsection*{Your turn}
\begin{itemize}
\raggedright
  \item \emph{Say it:} \emph{sulabhaḥ}
  \item \emph{Say it:} \emph{sulabhaḥ}, once more
  \item \emph{Say it:} say \emph{sulabhaḥ}
  \item \emph{Recall:} say the Sanskrit for that, then the Sanskrit for however, then say \emph{sulabhaḥ} again
\end{itemize}

\begin{tcolorbox}[breakable,colback=teal!4,colframe=teal!35!black,title={Before you move on}]
What does \sk{सुलभः} mean? (\textquotedblleft{}Easy to get, cheap\textquotedblright{}.) Say it once more.
\end{tcolorbox}
\section[\texorpdfstring{\sk{दुर्लभः} (durlabhaḥ)}{दुर्लभः (durlabhaḥ)}]{\sk{दुर्लभः} (durlabhaḥ) — hard to get, rare}

\label{lesson:SA-C242-durlabhah}



\begin{quote}
Before the new one: say the Sanskrit for however, then the Sanskrit for easy to get.
\end{quote}

\subsection*{You'll want to know: \sk{दुर्लभः}}
\textbf{\sk{दुर्लभः}} — \emph{durlabhaḥ} — \textquotedblleft{}hard to get, rare\textquotedblright{}.

\begin{grammarlens}[title={What you've built}]
One more word for this chapter.
\end{grammarlens}

\subsection*{Your turn}
\begin{itemize}
\raggedright
  \item \emph{Say it:} \emph{durlabhaḥ}
  \item \emph{Say it:} \emph{durlabhaḥ}, once more
  \item \emph{Say it:} say \emph{durlabhaḥ}
  \item \emph{Recall:} say the Sanskrit for however, then the Sanskrit for easy to get, then say \emph{durlabhaḥ} again
\end{itemize}

\begin{tcolorbox}[breakable,colback=teal!4,colframe=teal!35!black,title={Before you move on}]
What does \sk{दुर्लभः} mean? (\textquotedblleft{}Hard to get, rare\textquotedblright{}.) Say it once more.
\end{tcolorbox}
\section[\texorpdfstring{\sk{जटिलः} (jaṭilaḥ)}{जटिलः (jaṭilaḥ)}]{\sk{जटिलः} (jaṭilaḥ) — complicated}

\label{lesson:SA-C242-jatilah}



\begin{quote}
Before the new one: say the Sanskrit for easy to get, then the Sanskrit for hard to get.
\end{quote}

\subsection*{You'll want to know: \sk{जटिलः}}
\textbf{\sk{जटिलः}} — \emph{jaṭilaḥ} — \textquotedblleft{}complicated\textquotedblright{}.

\begin{grammarlens}[title={What you've built}]
One more word for this chapter.
\end{grammarlens}

\subsection*{Your turn}
\begin{itemize}
\raggedright
  \item \emph{Say it:} \emph{jaṭilaḥ}
  \item \emph{Say it:} \emph{jaṭilaḥ}, once more
  \item \emph{Say it:} say \emph{jaṭilaḥ}
  \item \emph{Recall:} say the Sanskrit for easy to get, then the Sanskrit for hard to get, then say \emph{jaṭilaḥ} again
\end{itemize}

\begin{tcolorbox}[breakable,colback=teal!4,colframe=teal!35!black,title={Before you move on}]
What does \sk{जटिलः} mean? (\textquotedblleft{}Complicated\textquotedblright{}.) Say it once more.
\end{tcolorbox}
\section[\texorpdfstring{\sk{सुकरः} (sukaraḥ)}{सुकरः (sukaraḥ)}]{\sk{सुकरः} (sukaraḥ) — easy to do}

\label{lesson:SA-C242-sukarah}



\begin{quote}
Before the new one: say the Sanskrit for hard to get, then the Sanskrit for complicated.
\end{quote}

\subsection*{You'll want to know: \sk{सुकरः}}
\textbf{\sk{सुकरः}} — \emph{sukaraḥ} — \textquotedblleft{}easy to do\textquotedblright{}.

\begin{grammarlens}[title={What you've built}]
One more word for this chapter.
\end{grammarlens}

\subsection*{Your turn}
\begin{itemize}
\raggedright
  \item \emph{Say it:} \emph{sukaraḥ}
  \item \emph{Say it:} \emph{sukaraḥ}, once more
  \item \emph{Say it:} say \emph{sukaraḥ}
  \item \emph{Recall:} say the Sanskrit for hard to get, then the Sanskrit for complicated, then say \emph{sukaraḥ} again
\end{itemize}

\begin{tcolorbox}[breakable,colback=teal!4,colframe=teal!35!black,title={Before you move on}]
What does \sk{सुकरः} mean? (\textquotedblleft{}Easy to do\textquotedblright{}.) Say it once more.
\end{tcolorbox}
\section[\texorpdfstring{\sk{दुष्करः} (duṣkaraḥ)}{दुष्करः (duṣkaraḥ)}]{\sk{दुष्करः} (duṣkaraḥ) — hard to do}

\label{lesson:SA-C242-duskarah}



\begin{quote}
Before the new one: say the Sanskrit for complicated, then the Sanskrit for easy to do.
\end{quote}

\subsection*{You'll want to know: \sk{दुष्करः}}
\textbf{\sk{दुष्करः}} — \emph{duṣkaraḥ} — \textquotedblleft{}hard to do\textquotedblright{}.

\begin{grammarlens}[title={What you've built}]
One more word for this chapter.
\end{grammarlens}

\subsection*{Your turn}
\begin{itemize}
\raggedright
  \item \emph{Say it:} \emph{duṣkaraḥ}
  \item \emph{Say it:} \emph{duṣkaraḥ}, once more
  \item \emph{Say it:} say \emph{duṣkaraḥ}
  \item \emph{Recall:} say the Sanskrit for complicated, then the Sanskrit for easy to do, then say \emph{duṣkaraḥ} again
\end{itemize}

\begin{tcolorbox}[breakable,colback=teal!4,colframe=teal!35!black,title={Before you move on}]
What does \sk{दुष्करः} mean? (\textquotedblleft{}Hard to do\textquotedblright{}.) Say it once more.
\end{tcolorbox}
